% SOURCE_STATEMENT: "Let $V$ be a possibly infinite-dimensional vector space and let $R:\\mathbb R\\to\\operatorname{GL}(V\\otimes V)$ satisfy the Yang–Baxter equation with spectral parameter. Can such families be studied by a unified algebraic or categorical theory, and how do analytic properties of $R(x)$ control the associated braided or quasi-braided categories?"
% Preparation and final self-audit: CONTRIBUTING.md prompts C and D; source/audit artifacts retained in the local physics_statements_2026_10_11 batch.
% !TeX program = lualatex
\documentclass[11pt]{article}
\usepackage[margin=25mm]{geometry}
\usepackage{fontspec}
\setmainfont{Latin Modern Roman}
\usepackage{amsmath,amssymb,mathtools}
\usepackage{unicode-math}
\setmathfont{Latin Modern Math}
\usepackage{xurl,hyperref}
\hypersetup{colorlinks=true,urlcolor=blue}
\setcounter{secnumdepth}{0}
\setlength{\parindent}{0pt}
\setlength{\parskip}{5pt}
\setlength{\emergencystretch}{3em}
\newcommand{\sref}[2]{\hyperlink{source:#1:#2}{[S#2]}}
\newcommand{\cataloguescope}[1]{\paragraph{Recorded target.}#1}
\newcommand{\R}{\mathbb R}
\newcommand{\C}{\mathbb C}
\newcommand{\N}{\mathbb N}
\newcommand{\Z}{\mathbb Z}
\newcommand{\Q}{\mathbb Q}
\newcommand{\F}{\mathbb F}
\newcommand{\eps}{\varepsilon}
\begin{document}
\section{Open Problems on Spectral-Parameter Yang–Baxter Solutions}\label{statement}

\subsection{Definitions and mathematical statement}

\paragraph{Shared notation.}$\mathbb N=\{1,2,\ldots\}$; the standard number fields have their usual meanings. All additional conventions are stated locally.


A spectral-parameter Yang--Baxter operator is an invertible map $R(x)$ on $V\otimes V$. Set $R_1(x)=R(x)\otimes I$ and $R_2(x)=I\otimes R(x)$. In the report's checked convention it satisfies, for every real $x,y$,
\[
R_1(x)R_2(x+y)R_1(y)=R_2(y)R_1(x+y)R_2(x).
\]
A braided category packages compatible exchanges; a quasi-braided variant allows the source's additional associativity data.

\paragraph{Catalogue wording retained for traceability.}
Let $V$ be a possibly infinite-dimensional vector space and let $R:\mathbb R\to\operatorname{GL}(V\otimes V)$ satisfy the Yang–Baxter equation with spectral parameter. Can such families be studied by a unified algebraic or categorical theory, and how do analytic properties of $R(x)$ control the associated braided or quasi-braided categories?

\paragraph{Scope and source qualification.}For infinite-dimensional $V$, domains and the topology of invertibility must be specified. Different checked and unchecked $R$ conventions give different written equations, so the report's convention is retained.
\subsection{Short English statement}
Develop a common theory of spectral-parameter Yang--Baxter families and relate their analytic behavior to braided or quasi-braided categories.
\subsection{Sources}
\begin{itemize}
\item[\hypertarget{source:statement:1}{S1}] ULAM.AI and the UnsolvedMath Contributors, UnsolvedMath, catalogue entry. Catalogue curation under CC BY 4.0; definitions and English exposition added locally.
\url{https://huggingface.co/datasets/ulamai/UnsolvedMath}.
\item[\hypertarget{source:statement:2}{S2}] Mini-Workshop: The Yang–Baxter Equation and Representations of Braid Groups (2025), Oberwolfach Reports, DOI 10.4171/owr/2025/49, pp. 2633--2672. Original report and mathematical context. Incorporated report material remains under the source CC BY-SA 4.0 terms unless otherwise noted.
\url{https://doi.org/10.4171/owr/2025/49}.
\end{itemize}
\label{end:statement}
\end{document}
